\chapter{これ・それ・あれ — Demostrativos}
\unidadheader{03}{これ・それ・あれ}{Demostrativos}{Señalar y preguntar por objetos y lugares}

\begin{tcolorbox}[vocabbox, title={📌 Vocabulario Nuevo}]
\begin{tabularx}{\linewidth}{llX}
\toprule
\textbf{Japonés} & \textbf{Español} & \textbf{Tipo} \\
\midrule
\vocabitem{これ}{—}{esto (cerca del hablante)}{pron.}
\vocabitem{それ}{—}{eso (cerca del oyente)}{pron.}
\vocabitem{あれ}{—}{aquello (lejos de ambos)}{pron.}
\vocabitem{ここ}{—}{aquí}{pron. lugar}
\vocabitem{そこ}{—}{ahí}{pron. lugar}
\vocabitem{あそこ}{—}{allá}{pron. lugar}
\vocabitem{この〜}{—}{este/esta ~ (+ sust.)}{det.}
\vocabitem{どんな}{—}{qué tipo de ~}{det.}
\vocabitem{何}{なに・なん}{qué}{pron.}
\vocabitem{誰}{だれ}{quién}{pron.}
\bottomrule
\end{tabularx}
\end{tcolorbox}

\begin{tcolorbox}[phrasebox, title={🗺️ Sistema こそあど}]
\begin{tabular}{lllll}
\toprule
 & \textbf{こ (cerca)} & \textbf{そ (oyente)} & \textbf{あ (lejos)} & \textbf{ど (pregunta)} \\
\midrule
cosa & これ & それ & あれ & どれ \\
+sust. & この & その & あの & どの \\
lugar & ここ & そこ & あそこ & どこ \\
dirección & こちら & そちら & あちら & どちら \\
tipo & こんな & そんな & あんな & どんな \\
\bottomrule
\end{tabular}
\end{tcolorbox}

\subsection*{🗣️ Frases Clave}

\frase{これは\ruby{何}{なん}ですか？}{¿Qué es esto?}
\frase{それは〜さんのですか？}{¿Es eso de ~san?}
\frase{\ruby{誰}{だれ}のですか？}{¿De quién es?}
\frase{あそこはどこですか？}{¿Qué lugar es aquel de allá?}

\subsection*{⚙️ Gramática}

\patron{Demostrativo + sustantivo}
  {この/その/あの/どの + 名詞}
  {\ej{この\ruby{本}{ほん}はいくらですか？}{¿Cuánto cuesta este libro?}
  \ej{あの\ruby{人}{ひと}は\ruby{誰}{だれ}ですか？}{¿Quién es aquella persona?}}

\patron{Posesión con の}
  {A の B = B de A}
  {\ej{これは\ruby{私}{わたし}の\ruby{本}{ほん}です。}{Este es mi libro.}
  \ej{\ruby{田中}{たなか}さんの\ruby{名前}{なまえ}は…}{El nombre del Sr. Tanaka es…}}

\patron{Pregunta de precio}
  {〜はいくらですか？}
  {\ej{このかばんはいくらですか？}{¿Cuánto cuesta esta bolsa?}
  \ej{\ruby{二千円}{にせんえん}です。}{Son 2000 yen.}}

\begin{tcolorbox}[ejerciciobox, title={📝 Ejercicios Rápidos}]
\begin{enumerate}
  \item Señala algo lejano a ambos y pregunta qué es.
  \item Completa: \_\_\_は\ruby{私}{わたし}のかばんです。(Este es mi bolso)
  \item Traduce: \ruby{誰}{だれ}のノートですか？
\end{enumerate}
\vspace{0.4em}
\textbf{Respuestas:}
1) あれは\ruby{何}{なん}ですか？\quad
2) これ\quad
3) ¿De quién es este cuaderno?
\end{tcolorbox}

\begin{tcolorbox}[kanjibox, title={🀄 Kanjis de la Unidad}]
\begin{tabularx}{\linewidth}{cllXX}
\toprule
\textbf{Kanji} & \textbf{On} & \textbf{Kun} & \textbf{Significado} & \textbf{Vocab} \\
\midrule
\kanjirow{何}{か}{なに・なん}{qué / cuánto}{\ruby{何}{なに} / \ruby{何時}{なんじ}}
\kanjirow{誰}{すい}{だれ}{quién}{\ruby{誰}{だれ} / \ruby{誰か}{だれか}}
\kanjirow{物}{ぶつ・もつ}{もの}{cosa / objeto}{\ruby{物}{もの} / \ruby{買い物}{かいもの}}
\bottomrule
\end{tabularx}
\end{tcolorbox}
